\documentclass[11pt]{article}
\usepackage[margin=25mm]{geometry}
\usepackage{fontspec}
\setmainfont{Latin Modern Roman}
\usepackage{amsmath,amssymb,mathtools}
\usepackage{unicode-math}
\setmathfont{Latin Modern Math}
\usepackage{xurl,hyperref}
\hypersetup{colorlinks=true,urlcolor=blue}
\setcounter{secnumdepth}{0}
\setlength{\parindent}{0pt}
\setlength{\parskip}{5pt}
\setlength{\emergencystretch}{3em}
\newcommand{\sref}[2]{\hyperlink{source:#1:#2}{[S#2]}}
\newcommand{\cataloguescope}[1]{\paragraph{Recorded target.}#1}
\newcommand{\R}{\mathbb R}
\newcommand{\C}{\mathbb C}
\newcommand{\N}{\mathbb N}
\newcommand{\Z}{\mathbb Z}
\newcommand{\Q}{\mathbb Q}
\newcommand{\F}{\mathbb F}
\newcommand{\eps}{\varepsilon}
\begin{document}
\section{Ericksen Universal Deformations of Incompressible Solids}\label{statement}

\subsection{Definitions and mathematical statement}

\paragraph{Shared notation.}$\mathbb N=\{1,2,\ldots\}$; the standard number fields have their usual meanings. All additional conventions are stated locally.


For a deformation $y$, write $F=\nabla y$ and $B=FF^T$. Incompressibility requires $\det F=1$. The principal invariants are $I_1=\operatorname{tr}B$, $I_2=\tfrac12[(\operatorname{tr}B)^2-\operatorname{tr}(B^2)]$ and $I_3=\det B$. A universal deformation satisfies equilibrium for every constitutive law in the stated homogeneous isotropic incompressible class, with an appropriate pressure multiplier.

\paragraph{Catalogue wording retained for traceability.}
Classify all constant-principal-invariant universal deformations of homogeneous incompressible isotropic solids.

\paragraph{Scope and source qualification.}Constant principal invariants do not mean that $F$ itself is spatially constant. The target includes possible nonaffine deformations.
\subsection{Short English statement}
Classify every universal deformation whose principal invariants remain constant throughout the incompressible isotropic solid.
\subsection{Research attempt}
\paragraph{Model, date and investigated scope.}
OpenAI GPT-6 Astra (configured model \texttt{gpt-6-astra}), ultra reasoning; generated 11 October 2026. Derived necessary and sufficient pressure-integrability conditions in a smooth simply connected constant-invariant setting, without classifying compatible deformations.

\paragraph{Sources and literature check.}
Read the full supplied definition, catalogue wording, source qualifications and listed references for this question. No fresh status or novelty search was performed. The argument below proves its auxiliary statements directly and does not import an unverified external theorem.

\paragraph{A differential reduction at constant invariants.} Restrict this test to a smooth three-dimensional incompressible deformation with spatial left Cauchy--Green tensor \(B\), on a simply connected current domain. For the usual isotropic energy \(W(I_1,I_2)\), the Cauchy stress has the form \(T=-pI+2W_1B-2W_2B^{-1}\), where \(W_i\) are derivatives of \(W\); an isotropic term is absorbed in the pressure \(p\). If both invariants are constant, \(W_1,W_2\) are constants in space. Thus the equilibrium equation becomes
\[
 \nabla p=2W_1\operatorname{div}B-2W_2\operatorname{div}B^{-1}.
\]
Taking the two admissible test energies \(W=I_1\) and \(W=I_2\) shows that universality requires both vector fields \(\operatorname{div}B\) and \(\operatorname{div}B^{-1}\) to be gradients.
\paragraph{Converse within this smooth setting.} If \(\operatorname{div}B=\nabla q\) and \(\operatorname{div}B^{-1}=\nabla r\), then choosing \(p=2W_1q-2W_2r\) proves equilibrium for every such \(W\). Equivalently, in a simply connected domain the two curls must vanish: integrating a curl-free field along paths defines a potential because integrals around contractible loops vanish by Stokes' formula. Affine maps \(y(X)=FX+c\) with \(\det F=1\) pass the test because both divergences vanish.
\paragraph{Why classification remains.} The reduction does not classify compatible fields \(B=\nabla y(\nabla y)^T\) with constant eigenvalues. Rotating their eigendirections arbitrarily generally destroys deformation compatibility or the two gradient conditions. On multiply connected domains vanishing curl also needs zero periods. The full source target includes nonaffine deformations and appropriate domains; solving this nonlinear compatibility system is the first unproved step, so the differential reformulation is not presented as a new classification.

\paragraph{Verification and adversarial audit.}
Checked both constitutive test energies, pressure signs, the affine case and the additional period obstruction on nonsimply connected domains. Model self-audit only; no independent human mathematical verification. Applied contribution prompt A before drafting and B afterwards. Checked both constitutive test energies, pressure signs, the affine case and the additional period obstruction on nonsimply connected domains. First remaining gap: Classify the compatible nonaffine deformation fields satisfying the two gradient conditions, including topology and source domain hypotheses. The inherited source status is preserved; original-target completion remains zero.

Coordinator checked the constant-invariant stress reduction, test-energy necessity and gradient-potential sufficiency on the stated domain. This additional check was model-assisted, not independent human verification.

\paragraph{Tools, process and contributions.}
Python 3 and local file tools used to read definitions and serialize manuscripts. Calculations are symbolic reasoning; no numerical experiment or theorem prover. Contribution prompt A was followed by prompt B. GPT-6 Astra with ultra reasoning drafted the research and performed its stated self-audit; any additional coordinator contribution is identified in the verification and audit. A GPT-6-based Codex coordinator assembled the manuscript and checked repository integration; its exact variant and reasoning configuration were not exposed. No human mathematical contribution, independent human proof check or formal proof certification is supplied. Complete derivations and the consulted source references are included here; private reasoning transcripts are not published.

\paragraph{Final status and first remaining gap.}
\textbf{UNRESOLVED}. Classify the compatible nonaffine deformation fields satisfying the two gradient conditions, including topology and source domain hypotheses.

\section{Completion Estimate}
\noindent\textbf{COMPLETION: 0\%}
This is an elementary auxiliary calculation or an explicit test of a restricted strategy. No new portion of the full source target is established.

\subsection{Sources}
\begin{itemize}
\item[\hypertarget{source:statement:1}{S1}] ULAM.AI and the UnsolvedMath Contributors, UnsolvedMath, catalogue entry. Catalogue curation under CC BY 4.0; definitions and English exposition added locally.
\url{https://huggingface.co/datasets/ulamai/UnsolvedMath}.
\item[\hypertarget{source:statement:2}{S2}] Mechanics of Materials: Multiscale Design of Advanced Materials and Structures (2024), Oberwolfach Reports, DOI 10.4171/owr/2024/11, pp. 657--730. Original report and mathematical context. Incorporated report material remains under the source CC BY-SA 4.0 terms unless otherwise noted.
\url{https://doi.org/10.4171/owr/2024/11}.
\end{itemize}
\label{end:statement}
\end{document}
